\documentclass{article}
\usepackage[UTF8]{ctex}
\usepackage{amsmath,amssymb,amsfonts}
\usepackage{enumitem}
\usepackage{geometry}
\geometry{a4paper, margin=1in}
% === 错题本排版优化（不参与题目标识） ===
\setlength{\parindent}{0pt}
\setlength{\parskip}{0.35em plus 0.10em minus 0.08em}
\setlength{\emergencystretch}{3em}
\setlist{itemsep=0.25em, topsep=0.40em, parsep=0pt}
\allowdisplaybreaks[2]
% === 错题本排版优化结束 ===
\begin{document}
\section*{厦门大学2024年数学分析真题}
\section*{113．厦门大学2024年数学分析试卷}

\subsection{题目1}
1．（15 分）求极限 \(\lim _{n \rightarrow \infty} \frac{1 \cdot 1!+2 \cdot 2!+\cdots+n \cdot n!}{(n+1)!}\) ．

\subsection{题目2}
2．（20 分）函数 \(f(x)\) 在 \([0,1]\) 上有二阶连续导数，且 \(f(0)=f(1)=0\) ，当 \(x \in(0,1)\) 时，\(f(x) \neq 0\) ．证明：

\[
\int_{0}^{1}\left|\frac{f^{\prime \prime}(x)}{f(x)}\right| \mathrm{d} x \geq 4
\]

\subsection{题目3}
3．（15 分）已知无穷积分 \(\int_{a}^{+\infty} f(x) \mathrm{d} x\) 收敛，且 \(f(x)\) 单调，证明： \(\lim _{x \rightarrow+\infty} x f(x)=0\) ．

\subsection{题目4}
4．（20 分）求二元函数 \(f(x, y)=x^{2}+y^{2}-12 x+16 y\) 在 \(D=\left\{(x, y) \mid x^{2}+y^{2} \leq 25\right\}\) 上的最值．

\subsection{题目5}
5．（20 分）有 \(K\) 是 \(n\) 维欧氏空间 \(\mathbb{R}^{n}\) 上的有界非空闭子集，且 \(K \subseteq \bigcup_{k=1}^{\infty} U_{k}\) ，其中 \(U_{k}\) 为开集．证明：存在 \(\varepsilon>0\) ，对于任意的 \(x \in K\) ，以 \(x\) 为中心 \(\varepsilon\) 为半径的开球 \(B(x, \varepsilon)\) 包含于某个 \(U_{k}\) 中．

\subsection{题目6}
6．（20 分）解答如下问题：

\begin{itemize}
  \item[（1）] 将 \(f(x)=\left\{\begin{array}{ll}-\frac{\pi}{4}, & x \in[-\pi, 0) ; \\ \frac{\pi}{4}, & x \in[0, \pi) .\end{array}\right.\) 展开成傅里叶级数．
  \item[（2）] 求 \(\sum_{n=1}^{\infty} \frac{1}{(2 n-1)^{2}}\) ．
\end{itemize}

\subsection{题目7}
7．（20 分）设 \(f(x)\) 有连续导数，且 \(f(0)=0\) ，求极限

\[
\lim _{t \rightarrow 0} \frac{1}{\pi t^{4}} \iiint_{x^{2}+y^{2}+z^{2} \leq t^{2}} f\left(\sqrt{x^{2}+y^{2}+z^{2}}\right) \mathrm{d} x \mathrm{~d} y \mathrm{~d} z .
\]

\subsection{题目8}
8．（20 分）计算曲线积分

\[
\int_{L} x \ln \left(x^{2}+y^{2}-1\right) \mathrm{d} x+y \ln \left(x^{2}+y^{2}-1\right) \mathrm{d} y .
\]

其中 \(L\) 为从 \((2,0)\) 到 \((0,2)\) 的光滑曲线．
\end{document}
